\textbf{COLUNI 2025 -- Questão 37}\par\medskip As décadas de 1920, 1930 e 1940 foram marcadas pela ascensão no nazismo alemão e do fascismo italiano na Europa. Hitler e Mussolini se tornaram lideranças políticas em seus países, perseguiram, mataram e exilaram seus inimigos. Em 1944, penúltimo ano da Segunda Guerra Mundial, o químico italiano, Primo Levi, foi feito prisioneiro político pelo regime fascista e levado para um campo de concentração na Itália. Logo depois, foi enviado para o campo de Auschwitz, na Polônia. Sendo um sobrevivente do holocausto, Levi legou à humanidade um dos mais valiosos testemunhos sobre o horror do período. Em uma de suas obras, ele escreveu:

\par

\begin{quote}
Imagine-se, agora, um homem privado não apenas dos seres queridos, mas de sua casa, seus hábitos, sua roupa, tudo, enfim, rigorosamente tudo que possuía; ele será um ser vazio, reduzido a puro sofrimento e carência, esquecido de dignidade e discernimento... facilmente se decidirá sobre sua vida e sua morte, sem qualquer sentimento de afinidade humana. Ficará claro, então, o duplo significado da expressão “Campo de extermínio”... aprendi que sou um Haftling. Meu nome é 174.517; fomos batizados, levaremos até a morte essa marca tatuada no braço esquerdo. A operação foi pouco dolorosa e extraordinariamente rápida: colocaram-nos numa fila e, um por um, conforme a ordem alfabética dos nossos nomes, passamos por um hábil funcionário, munido de uma espécie de punção com uma agulha minúscula... só “mostrando o número” recebe-se o pão e a sopa... durante muitos dias, quando o hábito da vida em liberdade me levava a olhar a hora no relógio, no pulso aparecia-me, ironicamente, meu novo nome, esse número tatuado em marcas azuladas sob a pele.
\end{quote}
\begin{flushright}LEVI, Primo. \textbf{Isto é um homem?} Rio de Janeiro: Rocco, 1988, p. 25.\end{flushright}

\par

A partir do trecho da obra de Primo Levi e do contexto de ascensão e queda do nazismo e do fascismo na Europa, marque a afirmativa CORRETA:\begin{enumerate}
\item Hitler e Mussolini chegaram ao poder por via de golpes de Estado, sendo Mussolini ovacionado na “Marcha
        sobre Roma”, de 1922, que derrubou o rei Victor Emanuel III, e Hitler, apoiado pelos “camisas pardas” de
        Munique, na deposição do presidente e chanceler Paul Von Hindenburg.
\item Hitler, em seu livro \textit{Mein Kampf} (“Minha Luta”), de 1925, apresenta um projeto de raça, de nação e de
        educação, a fim de educar o povo alemão na arte do conhecimento, no respeito às diferenças e no apoio
        irrestrito ao Estado.
\item Nos campos de concentração, como em Auschwitz, não havia diferença no tratamento dado aos
        prisioneiros, sendo todos tratados com a mesma violência, afinal, como fica claro no texto de Levi, eram
        todos judeus perigosos e delinquentes.
\item O racismo científico, usado por Hitler para justificar seu ideal de purificação, foi um instrumento poderoso
        de coisificação do homem nos campos de concentração, fazendo dos prisioneiros objetos sem vontade,
        sem identidade e sem dignidade.\end{enumerate}
